%
% $Log: graphs.tex,v $
% Revision 1.5  2002/10/21 15:26:44  mjm
% to ohio, revised
%
% Revision 1.4  2002/01/08 18:53:18  mjm
% made pdf compatible
%
% Revision 1.3  2001/08/15 18:07:03  mjm
% APPM preprint 466
%
% Revision 1.2  1999/08/09 19:55:27  mjm
% beta release
%
% Revision 1.1  1999/06/10 23:30:48  mjm
% Initial revision
%
%

%\begin{probdescript}{Drawing Graphs}
\handout{Graphs}

Graphs are an important visual aid in mathematics and it is important to
learn how to draw them properly.  For a graph to be useful it must display
the important and interesting information about a function and have all
relevant points clearly labeled.  

When drawing graphs, it is helpful to
ask yourself:
\begin{itemize}
\item Have I graphed what is being asked?  
\item Have I put a title on the  graph and labeled the axes?  
\item Is the scale I have chosen for my graph appropriate?
\item Have I illustrated the important features?
\item Have I included enough detail?
%  For example, labels on extreme
%  points, points of inflection, and $x$- and $y$-intercepts.
\end{itemize}

You can often determine how much detail to include in your graph from the
question.  For example, if you are finding the extreme points of a
function, then you will probably be expected to label these points on your
graph.  If you are only sketching the graph of a function then less detail
is required.  The main purpose of a sketch is to get some idea of the
general shape and behavior of the function.  In this case, it is not
usually necessary to label extreme points and intercepts.  However, you
should always remember to label your axes and put a title on your graph,
even if it is only a sketch.  Use a ruler if necessary to draw
straight lines. If 
you are drawing more than one curve on the same set of axes, take
particular care to make sure that each curve is clearly labeled.

%\end{probdescript}
%\clearpage

%\begin{probexamples}
%\begin{examplelist}
%\begin{example}

%\noindent{\bf Examples:}
\subsection*{Examples:}

  Graph the curve $y= x^3 - 4x$.

\begin{goodexample}

\setlength{\unitlength}{1cm}
\begin{center}
The curve $y=x^3-4x$.\\
\smallskip
\begin{picture}(10,8)(-5,-4)

% make axes

\put(0,-4){\vector(0,1){8}}
\put(-5,0){\vector(1,0){10}}
\thinlines
\multiput(-4,-0.1)(1,0){9}{\line(0,1){0.2}}

\multiput(-0.1,-3.5)(0,0.5){15}{\line(1,0){0.2}}

\put(5.2,0.2){$x$}
\put(0.2,3.8){$y$}

\put(-4.4,-0.4){$-4$}
\put(-3.4,-0.4){$-3$}
\put(-2.4,-0.4){$-2$}
\put(-1.4,-0.4){$-1$}

\put(0.9,-0.4){$1$}
\put(1.9,-0.4){$2$}
\put(2.9,-0.4){$3$}
\put(3.9,-0.4){$4$}

\put(-0.7,-2.6){$-5$}
\put(-0.4,2.4){$5$}
% add in the curve x^3+4x

%% FAKE but looks same

\qbezier(-2.6,-3.6)(-1.7,4.35)(0,0)
\qbezier(0,0)(1.7,-4.35)(2.6,3.6)

%\path(-2.6,-3.588)
%(-2.5,-2.8125)
%(-2.4,-2.112)
%(-2.3,-1.4835)
%(-2.2,-0.924)
%(-2.1,-0.4305)
%(-2,0)
%(-1.9,0.3705)
%(-1.8,0.684)
%(-1.7,0.9435)
%(-1.6,1.152)
%(-1.5,1.3125)
%(-1.4,1.428)
%(-1.3,1.5015)
%(-1.2,1.536)
%(-1.1,1.5345)
%(-1,1.5)
%(-0.9,1.4355)
%(-0.8,1.344)
%(-0.7,1.2285)
%(-0.6,1.092)
%(-0.5,0.9375)
%(-0.4,0.768)
%(-0.3,0.5865)
%(-0.2,0.396)
%(-0.1,0.1995)
%(0,0)
%(0.1,-0.1995)
%(0.2,-0.396)
%(0.3,-0.5865)
%(0.4,-0.768)
%(0.5,-0.9375)
%(0.6,-1.092)
%(0.7,-1.2285)
%(0.8,-1.344)
%(0.9,-1.4355)
%(1,-1.5)
%(1.1,-1.5345)
%(1.2,-1.536)
%(1.3,-1.5015)
%(1.4,-1.428)
%(1.5,-1.3125)
%(1.6,-1.152)
%(1.7,-0.9435)
%(1.8,-0.684)
%(1.9,-0.3705)
%(2,0)
%(2.1,0.4305)
%(2.2,0.924)
%(2.3,1.4835)
%(2.4,2.112)
%(2.5,2.8125)
%(2.6,3.588)

% add the labels

%\put(-2,0){\circle*{0.1}}
\put(-2.08,-0.08){\mbox{$\bullet$}}
\put(-3.2,0.3){$(-2,0)$}

%\put(-1.1,1.5396){\circle*{0.1}}
\put(-1.18,1.45){\mbox{$\bullet$}}
\put(-2.5,2){$(\frac{-2}{\sqrt{3}}, \frac{16}{3\sqrt{3}})$}


%\put(0,0){\circle*{0.1}}
\put(-0.08,-0.08){\mbox{$\bullet$}}
\put(0.1,0.5){$(0,0)$}

%\put(1.1,-1.5396){\circle*{0.1}}
\put(1.02,-1.62){\mbox{$\bullet$}}
\put(0.5,-2){$(\frac{2}{\sqrt{3}}, \frac{-16}{3\sqrt{3}})$}

%\put(2,0){\circle*{0.1}}
\put(1.92,-0.08){\mbox{$\bullet$}}
\put(2.4,0.5){$(2,0)$}



\end{picture}
\end{center}
\noindent\fbox{\parbox{\textwidth}{\rm The curve,
axes, extreme points, and $x$- and $y$-intercepts are clearly labeled.
These are the `points of interest' for this function.}}

\end{goodexample}

\secondpage

\begin{badexample}

\setlength{\unitlength}{0.7cm}
\begin{center}
\begin{picture}(10,8)(-5,-4)

% make axes

\put(0,-4){\vector(0,1){8}}
\put(-5,0){\vector(1,0){10}}
\thinlines
\multiput(-4,-0.1)(1,0){9}{\line(0,1){0.2}}

%\multiput(-0.1,-3.5)(0,0.5){15}{\line(1,0){0.2}}

\put(5.2,0.2){$x$}
\put(0.2,3.8){$y$}

\put(-4.4,-0.4){$-4$}
\put(-3.4,-0.4){$-3$}
\put(-2.4,-0.4){$-2$}
\put(-1.4,-0.4){$-1$}

\put(0.9,-0.4){$1$}
\put(1.9,-0.4){$2$}
\put(2.9,-0.4){$3$}
\put(3.9,-0.4){$4$}

% add in the curve x^3+4x

% FAKE 
\qbezier(-2.6,-3.6)(-1.7,4.35)(0,0)
\qbezier(0,0)(1.7,-4.35)(2.6,3.6)

\end{picture}
\end{center}


\noindent\fbox{\parbox{\textwidth}{\rm 
Although this second solution has labeled axes, the graph has no title
so it is not clear to someone else what they are viewing.  There are
also no labels on the extreme points or intercepts and no scale is given
on the $y$-axis.  Although this is a bad example of a solution to this
question it  would be  sufficient if only a rough sketch
was required.
}}
\end{badexample}


Choose a suitable scale so that all of the interesting features of the
graph are included.  For example, when graphing our function
\begin{displaymath}
  y= x^3 - 4x
\end{displaymath}
a suitable scale meant including all of the extreme points and $x$- and
$y$-intercepts.  
Be especially careful about the scale when using a graphing calculator or
computer. They do not have any idea which points are interesting, so they
will often give inappropriate scales.
The following graphs
give examples of bad choices of scale for our problem. 

\begin{badexample}


\setlength{\unitlength}{0.55cm}

\begin{center}
\begin{picture}(10,8)(-5,-4)
% make axes
\put(0,-4){\vector(0,1){8}}
\put(-5,0){\vector(1,0){10}}
\thinlines
%\multiput(-4,-0.1)(1,0){9}{\line(0,1){0.2}}
%\multiput(-0.1,-3.5)(0,0.5){15}{\line(1,0){0.2}}
\put(5.2,0.2){$x$}
\put(0.2,3.8){$y$}

\put(-4.6,-0.6){$-40$}
\put(3.9,-0.6){$40$}
\put(-2,3){$60000$}
\put(-2.5,-4){$-60000$}

\qbezier(-5,-4)(-3,0)(0,0)
\qbezier(5,4)(3,0)(0,0)

\end{picture}
%%%%%%
\hfill
%%%%%%
\begin{picture}(10,8)(-5,-4)

\qbezier(-5,4)(-3,3.5)(0,0)
\qbezier(5,-4)(3,-3.5)(0,0)

% make axes

\put(0,-4){\vector(0,1){8}}
\put(-5,0){\vector(1,0){10}}
\thinlines

\put(5.2,0.2){$x$}
\put(0.2,3.8){$y$}

\put(-4.4,-0.6){$-1$}
\put(3.9,-0.6){$1$}

\put(-1,3){$1$}
\put(-1.5,-4){$-1$}

\end{picture}

\end{center}


\noindent\fbox{\parbox{\textwidth}{\rm 
The graph on the left uses
too large a scale so the extreme points are hard to see and will be
difficult to label.  The scale used for the graph on the right is too 
small as it fails to include two of the $x$-intercepts.
}}
\end{badexample}

Also make sure that what you are graphing makes sense.  If you are drawing
the graph of a function that represents the distance of a space shuttle
from the surface of the earth against time, it makes no sense to include
negative distances, as this is unphysical.  Similar reasoning applies to
graphs representing populations, areas, and volumes.

Finally, being able to draw and sketch graphs well can sometimes help
while trying to answer a question, even if a graph is not required in the
solution.  For more difficult problems in particular, being able to sketch
a function to identify its behavior often provides more insight than
looking at the equation alone.
